% !TEX spellcheck = en_US



% ~~~~~~~~~~~~~~~~~~~~~~~~~~~~~~~~~~~~~~~ V
% (NC) 2026 Haluk Bingol
% github.com/halukbingol/LaTeX-Templates
%
% Licensees may copy, distribute, display, and perform the work and make derivative works 
% and remixes based on it only for non-commercial purposes. 
% ~~~~~~~~~~~~~~~~~~~~~~~~~~~~~~~~~~~~~~~ A


% =====================================================================
%  Code + Output presentation template
%  ---------------------------------------------------------------------
%  The macros live in the shared libraries hblibCodeoutput.sty and
%  hblibPresentation.sty (see their headers). This file loads them and
%  shows examples.
%  run_examples.sh runs codeoutput/<name>.<ext> and writes codeoutput/<name>.txt.
% =====================================================================

\documentclass[aspectratio=169,10pt,compress]{beamer}

% ~~~~~~~~~~~~~~~~~~~~~~~~~~~~~~~~~~~~~~~~~~~~~~~~~~~~~~~~~~~~~~ V beamer
	\useoutertheme{split}
	\useinnertheme{rounded}
	\usecolortheme{whale}
	\usecolortheme{orchid}
	\setbeamertemplate{blocks}[rounded]
	\setbeamertemplate{navigation symbols}{}
	% "i / N" on every frame comes from hblibPresentation.sty (\insertshorttitle)
% ~~~~~~~~~~~~~~~~~~~~~~~~~~~~~~~~~~~~~~~~~~~~~~~~~~~~~~~~~~~~~~ A beamer

% ~~~~~~~~~~~~~~~~~~~~~~~~~~~~~~~~~~~~~~~~~~~~~~~~~~~~~~~~~~~~~~~~~ V hblib
	\usepackage{../../hblib/hblibPresentation} % lib presentation 
	\usepackage{../../hblib/hblibCodeoutput} % lib code listing
% ~~~~~~~~~~~~~~~~~~~~~~~~~~~~~~~~~~~~~~~~~~~~~~~~~~~~~~~~~~~~~~~~~ A hblib


% xcolor, array, listings, tikz and algorithm2e come from the hblib libraries
\usepackage[T1]{fontenc}
\usepackage{lmodern}
\usepackage{booktabs}

% ---------------------------------------------------------------------
%  algorithm2e set-up (the package is loaded by hblibPresentation)
% ---------------------------------------------------------------------
\SetAlCapFnt{\small}
\SetAlCapNameFnt{\small}
\DontPrintSemicolon
\SetKwProg{Fn}{function}{}{end}
\SetKwFunction{FDijkstra}{Dijkstra}
\SetKwFunction{FExtractMin}{ExtractMin}
\SetKwFunction{FDecreaseKey}{DecreaseKey}
\SetKw{KwNil}{nil}

\newcolumntype{L}[1]{>{\raggedright\arraybackslash}p{#1}}

\definecolor{gitred}{RGB}{240,80,50}


%\newcommand{\cmd}[1]{\texttt{\color{gitred}#1}}

% ---------------------------------------------------------------------
\title{hb Cis \LaTeX\ Template}
\subtitle{A Template That Reads Both from Files}
\author[Bingol]{Haluk O. Bingol}

\institute[FCIS]{
    Faculty of Computer and Information Sciences\\
    Yeditepe University
}
\date{
	{\tiny \hbTimeStamp}
}

\begin{document}




% ~~~~~~~~~~~~~~~~~~~~~~~~~~~~~~~~~~~~~~ V frm
\begin{frame}[t,fragile]
  \titlepage
\end{frame}
% ~~~~~~~~~~~~~~~~~~~~~~~~~~~~~~~~~~~~~~ A frm




% ~~~~~~~~~~~~~~~~~~~~~~~~~~~~~~~~~~~~~~ V frm
\begin{frame}[t,fragile] \frametitle{Outline}
  \tableofcontents[hideallsubsections]
\end{frame}
% ~~~~~~~~~~~~~~~~~~~~~~~~~~~~~~~~~~~~~~ A frm




% ~~~~~~~~~~~~~~~~~~~~~~~~~~~~~~~~~~~~~~ V
%\hSection[General]{General}{}{}
\hSection[General]{General}
% ~~~~~~~~~~~~~~~~~~~~~~~~~~~~~~~~~~~~~~ A




% ~~~~~~~~~~~~~~~~~~~~~~~~~~~~~~~~~~~~~~ V
\subsection[hStuff]{\hCode{hStuff}}
% ~~~~~~~~~~~~~~~~~~~~~~~~~~~~~~~~~~~~~~ A




% ~~~~~~~~~~~~~~~~~~~~~~~~~~~~~~~~~~~~~~ V
\subsection[menukeys]{\hCode{menukeys} package}
% ~~~~~~~~~~~~~~~~~~~~~~~~~~~~~~~~~~~~~~ A




% ~~~~~~~~~~~~~~~~~~~~~~~~~~~~~~~~~~~~~~ V frm
\begin{frame}[t,fragile] \frametitle{\hCode{menukeys} package- Keys}
	
	Use \hCode{keys} to represent key strokes on the keyboard.
	One key \keys{X} or  key combinations such as
	\keys{Ctrl+X} or 
	\keys{Ctrl+Shift+F}.

	Common key combinations:
	
	%: +++++++++++++++++++++++++++++++++++++++ V -Tbl. cut-paste
	\begin{center}
	\sffamily\small
	\begin{tabular}{|l |l |}
		\hline
		\textbf{function} &\textbf{key combination}\\
		\hline
		cut &\keys{Ctrl+X}\\
		paste &\keys{Ctrl+V}\\
		help in eclipse &\keys{Ctrl+Space}\\
		\hline
	\end{tabular}
	\end{center}
	% +++++++++++++++++++++++++++++++++++++++ A
	
	Special keys

	%: +++++++++++++++++++++++++++++++++++++++ V -Tbl. special keys
	\begin{center}
	\sffamily\small
	\begin{tabular}{|l |l  |l |}
		\hline
		\textbf{function} &\textbf{key combination}&\textbf{key symbol}\\
		\hline
		Control &\keys{Ctrl} &\keys{\ctrl}\\
		Command &\keys{Cmd} &\keys{\cmd}\\
		Escape &\keys{Esc} &\keys{\esc}\\
		Shift &\keys{Shift} &\keys{\shift}\\
		Alt &\keys{Alt} &\keys{\Alt}\\
		Windows &\keys{Win} &none\\
		\hline
	\end{tabular}
	\end{center}
	% +++++++++++++++++++++++++++++++++++++++ A

\end{frame}
% ~~~~~~~~~~~~~~~~~~~~~~~~~~~~~~~~~~~~~~ A frm




% ~~~~~~~~~~~~~~~~~~~~~~~~~~~~~~~~~~~~~~ V frm
\begin{frame}[t,fragile] \frametitle{\hCode{menukeys} package- Menu}
	
	Use \hCode{menu} to describe menu sequences.
	
	\textbf{Example}
	In eclipse, source code is formatted with key combination
	\keys{Ctrl+Shift+F}.
	It can be done by menu, too:
	\begin{itemize}
		
		\item 
		In context menu 
		\menu{Source > Format}.
	
		\item 
		On top menu \menu{Source > Format}.
		
	\end{itemize}

	It may be multiple level deep
	\menu{Extras > Settings > Saving} or simple one level 
	\menu{Change save path}.
\end{frame}
% ~~~~~~~~~~~~~~~~~~~~~~~~~~~~~~~~~~~~~~ A frm




% ~~~~~~~~~~~~~~~~~~~~~~~~~~~~~~~~~~~~~~ V frm
\begin{frame}[t,fragile] \frametitle{\hCode{menukeys} package- Path}
	
	Use \hCode{\directory} to represent 
	a path such as\\
	\hCode{Macintosh HD/Users/Your Name/Documents} \\
	as\\
	\directory{Macintosh HD/Users/Your Name/Documents}.

\end{frame}
% ~~~~~~~~~~~~~~~~~~~~~~~~~~~~~~~~~~~~~~ A frm




% ~~~~~~~~~~~~~~~~~~~~~~~~~~~~~~~~~~~~~~ V frm
\begin{frame}[t,fragile] \frametitle{hCode}
	Use \hCode{your code here} to talk about the code.
%	aaa \hKeyboard{a} bbb
%	aaa \hKeyboard{Ctrl}+\hKeyboard{Cmd}+\hKeyboard{Shft} bbb

\end{frame}
% ~~~~~~~~~~~~~~~~~~~~~~~~~~~~~~~~~~~~~~ A frm




% ~~~~~~~~~~~~~~~~~~~~~~~~~~~~~~~~~~~~~~ V
\hSection[Pseudo code]{Pseudo code}
% ~~~~~~~~~~~~~~~~~~~~~~~~~~~~~~~~~~~~~~ A




% ~~~~~~~~~~~~~~~~~~~~~~~~~~~~~~~~~~~~~~ V
\subsection[Dijkstra'sSP]{Dijkstra's algorithm}
% ~~~~~~~~~~~~~~~~~~~~~~~~~~~~~~~~~~~~~~ A




% ~~~~~~~~~~~~~~~~~~~~~~~~~~~~~~~~~~~~~~ V frm
\begin{frame}[t,fragile] \frametitle{Pseudo code}
  \begin{columns}[T]
    \begin{column}{0.62\textwidth}
      \SetAlFnt{\scriptsize}


% +++++++++++++++++++++++++++++++++++++++ V alg
%: -Alg. 
\begin{algorithm}[H]
  \caption{Dijkstra's algorithm}
  \KwIn{graph $G = (V, E)$, weights $w(u, v) \ge 0$, source $s$}
  \KwOut{distances $d[\,]$ and predecessors $\pi[\,]$}
  \Fn{\FDijkstra{$G, w, s$}}{
    \ForEach{$v \in V$}{
      $d[v] \leftarrow \infty$\;
      $\pi[v] \leftarrow$ \KwNil\;
    }
    $d[s] \leftarrow 0$\;
    $Q \leftarrow$ priority queue of all $v \in V$, keyed by $d[v]$\;
    \While{$Q \neq \emptyset$}{
      $u \leftarrow$ \FExtractMin{$Q$}\tcp*{$d[u]$ is now final}
      \ForEach{edge $(u, v) \in E$}{
        \If(\tcp*[f]{relax $(u, v)$}){$d[u] + w(u, v) < d[v]$}{
          $d[v] \leftarrow d[u] + w(u, v)$\;
          $\pi[v] \leftarrow u$\;
          \FDecreaseKey{$Q, v, d[v]$}\;
        }
      }
    }
    \Return{$d, \pi$}\;
  }
\end{algorithm}
% +++++++++++++++++++++++++++++++++++++++ A alg

    \end{column}
    \begin{column}{0.35\textwidth}
      \begin{itemize}\small
        \item \textbf{Lines 2--6}: every vertex starts \emph{unreached}, except $s$
        \item \textbf{Line 8}: the vertex closest to $s$ among the unfinished ones
        \item \textbf{Lines 10--13}: \emph{relaxation}: a shorter way to $v$ through $u$?
        \item \texttt{DecreaseKey} moves $v$ forward in the queue
        \item $\infty + w = \infty$, so unreached vertices never improve others
      \end{itemize}
    \end{column}
  \end{columns}
\end{frame}
% ~~~~~~~~~~~~~~~~~~~~~~~~~~~~~~~~~~~~~~ A frm




% ~~~~~~~~~~~~~~~~~~~~~~~~~~~~~~~~~~~~~~ V
\section[Usage]{How to Use This Template}
% ~~~~~~~~~~~~~~~~~~~~~~~~~~~~~~~~~~~~~~ A




% ~~~~~~~~~~~~~~~~~~~~~~~~~~~~~~~~~~~~~~ V frm
\begin{frame}[t,fragile] \frametitle{Project layout and workflow}
  \begin{columns}[T]
    \begin{column}{0.42\textwidth}


% +++++++++++++++++++++++++++++++++++++++ V lst
%: -Lst.
\begin{lstlisting}[language={}, basicstyle=\ttfamily\footnotesize, frame=single,
                   backgroundcolor=\color{codebg}, rulecolor=\color{gray!40}]
zCisLaTeXTemplate/hblib/
|-- hblibCodeoutput.sty
`-- hblibPresentation.sty
course/project/
|-- hbCisLaTeXTemplate.tex
|-- run_examples.sh
|-- codeoutput/
|   |-- fibonacci.py
|   |-- fibonacci.txt
|   |-- factorial.c
|   |-- factorial.txt
|   |-- Hello.java
|   `-- Hello.txt
\end{lstlisting}
% +++++++++++++++++++++++++++++++++++++++ A lst


    \end{column}
    \begin{column}{0.54\textwidth}
      \begin{enumerate}\small
        \item Load the libraries: \verb|\usepackage{hblibCodeoutput}|, \verb|\usepackage{hblibPresentation}|
        \item Put each program in \texttt{codeoutput/}
        \item Run \cmd{sh run\_examples.sh}. It runs every program and saves what it prints
              (stdout \emph{and} stderr) to \texttt{codeoutput/<name>.txt}
        \item Compile the slides: \cmd{pdflatex hbCisLaTeXTemplate.tex}
      \end{enumerate}

      \begin{block}{One source of truth}
        \small Edit the program, re-run the script, recompile. The code and its output on
        the slides can never disagree, because neither is copied into the \texttt{.tex}.
      \end{block}
    \end{column}
  \end{columns}
\end{frame}
% ~~~~~~~~~~~~~~~~~~~~~~~~~~~~~~~~~~~~~~ A frm




% ~~~~~~~~~~~~~~~~~~~~~~~~~~~~~~~~~~~~~~ V frm
\begin{frame}[t,fragile] \frametitle{The macros}
  \footnotesize
  \begin{tabular}{@{}L{0.28\textwidth}L{0.66\textwidth}@{}}
    \toprule
    \textbf{Macro} & \textbf{Shows} \\
    \midrule
    \texttt{\textbackslash lstcode}        & code only \\
    \texttt{\textbackslash lstoutput}      & output only \\
    \texttt{\textbackslash lstcodedown}         & code, then output below \\
    \texttt{\textbackslash lstcoderight}     & code left, output right \\
    \midrule
    \multicolumn{2}{@{}l}{\emph{Code on the first slide, output added on the second:}} \\
    \texttt{\textbackslash lstcodedownstep}     & output below the code \\
    \texttt{\textbackslash lstcoderightstep} & output to the right of the code \\
    \texttt{\textbackslash lstcodeoverlaystep}  & output floats on top, upper-left corner at $(x,y)$ \\
    \bottomrule
  \end{tabular}

  \vspace{0.2em}


% +++++++++++++++++++++++++++++++++++++++ V lst
%: -Lst.
\begin{lstlisting}[language={[LaTeX]TeX}, basicstyle=\ttfamily\scriptsize, frame=single,
                   backgroundcolor=\color{codebg}, rulecolor=\color{gray!40}]
\lstcode           [code opts]{code path}{Language}
\lstoutput         [output opts]{output path}
\lstcodedown       [code opts][output opts]{code path}{Language}{output path}
\lstcoderight      [code width][code opts][output opts]{code path}{Language}{output path}
\lstcodedownstep   [code opts][output opts]{code path}{Language}{output path}
\lstcoderightstep  [code width][code opts][output opts]{code path}{Language}{output path}
\lstcodeoverlaystep[code opts][output opts][output width]{x,y}
                   {code path}{Language}{output path}
\end{lstlisting}
% +++++++++++++++++++++++++++++++++++++++ A lst


\end{frame}
% ~~~~~~~~~~~~~~~~~~~~~~~~~~~~~~~~~~~~~~ A frm




% ~~~~~~~~~~~~~~~~~~~~~~~~~~~~~~~~~~~~~~ V frm
\begin{frame}[t,fragile] \frametitle{Paths and options}
  \begin{itemize}\small
    \item Paths are \textbf{relative to the \texttt{.tex} file}, so code and output can live
          anywhere: \texttt{codeoutput/primes.py}, \texttt{../week3/sort.c}, \texttt{lab2/out.txt}
    \item The code and output paths are given separately, so one program can have
          several outputs (e.g.\ \texttt{codeoutput/run1.txt}, \texttt{codeoutput/run2.txt})
    \item Optional arguments take ordinary \texttt{listings} options:
  \end{itemize}


% +++++++++++++++++++++++++++++++++++++++ V lst
%: -Lst.
\begin{lstlisting}[language={[LaTeX]TeX}, basicstyle=\ttfamily\scriptsize, frame=single,
                   backgroundcolor=\color{codebg}, rulecolor=\color{gray!40}]
\lstcode[firstline=4, lastline=9, firstnumber=4]{codeoutput/factorial.c}{C}  % excerpt
\lstcodedown[][lastline=5]{codeoutput/primes.py}{Python}{codeoutput/primes.txt}  % 5 output lines
\lstcoderight[0.6][basicstyle=\ttfamily\scriptsize]
   {codeoutput/git-demo.sh}{bash}{codeoutput/git-demo.txt}                       % smaller code
\lstcodeoverlaystep[][][0.4\textwidth]{8.6cm,2.3cm}
   {codeoutput/primes.py}{Python}{codeoutput/primes.txt}                         % at (8.6,2.3)
\end{lstlisting}
% +++++++++++++++++++++++++++++++++++++++ A lst


\end{frame}
% ~~~~~~~~~~~~~~~~~~~~~~~~~~~~~~~~~~~~~~ A frm




% ~~~~~~~~~~~~~~~~~~~~~~~~~~~~~~~~~~~~~~ V frm
\begin{frame}[t,fragile] \frametitle{Sharing the library}
  All macros live in \texttt{hblibCodeoutput.sty}. Load it in any beamer presentation:


% +++++++++++++++++++++++++++++++++++++++ V lst
%: -Lst.
\begin{lstlisting}[
		language={[LaTeX]TeX}, basicstyle=\ttfamily\scriptsize, frame=single,
		backgroundcolor=\color{codebg}, rulecolor=\color{gray!40}
		]
	\usepackage{hblibCodeoutput}
	\definecolor{codeframe}{RGB}{120,40,120}   % optional: recolor after loading
\end{lstlisting}
% +++++++++++++++++++++++++++++++++++++++ A lst


  \begin{columns}[T]
    \begin{column}{0.52\textwidth}
      \textbf{Where to put the file}
      \begin{itemize}\small
        \item \textbf{One project:} next to the \texttt{.tex} file
        \item \textbf{All projects:} in your personal texmf tree
              \begin{itemize}\footnotesize
                \item Windows (MiKTeX): add a folder in \emph{MiKTeX Console, Settings, Directories}
                      and put the files in its \texttt{tex/latex/hblib/}
                \item macOS: \texttt{\textasciitilde/Library/texmf/tex/latex/hblib/}
              \end{itemize}
        \item Check with \cmd{kpsewhich hblibCodeoutput.sty}
      \end{itemize}
    \end{column}
    \begin{column}{0.44\textwidth}
      \textbf{Colors you can change}
      \begin{itemize}\small
        \item code: \texttt{codebg}, \texttt{codeframe}, \texttt{codekw}, \texttt{codestr}, \texttt{codecmt}
        \item output: \texttt{termbg}, \texttt{termfg}, \texttt{termframe}
      \end{itemize}
      \small The full macro list is in the header of \texttt{hblibCodeoutput.sty}.
    \end{column}
  \end{columns}
\end{frame}
% ~~~~~~~~~~~~~~~~~~~~~~~~~~~~~~~~~~~~~~ A frm

% ~~~~~~~~~~~~~~~~~~~~~~~~~~~~~~~~~~~~~~ V
\hSection[CodeLocation]{Code location}{}{}
% ~~~~~~~~~~~~~~~~~~~~~~~~~~~~~~~~~~~~~~ A


% ~~~~~~~~~~~~~~~~~~~~~~~~~~~~~~~~~~~~~~ V
\subsection[down]{Output at below}
% ~~~~~~~~~~~~~~~~~~~~~~~~~~~~~~~~~~~~~~ A




% ~~~~~~~~~~~~~~~~~~~~~~~~~~~~~~~~~~~~~~ V frm
\begin{frame}[t,fragile] \frametitle{Code, then output: \texttt{\textbackslash lstcodedown}}
  \lstcodedown{codeoutput/fibonacci.py}{Python}{codeoutput/fibonacci.txt}
\end{frame}
% ~~~~~~~~~~~~~~~~~~~~~~~~~~~~~~~~~~~~~~ A frm





% ~~~~~~~~~~~~~~~~~~~~~~~~~~~~~~~~~~~~~~ V
\subsection[right]{Output at right}
% ~~~~~~~~~~~~~~~~~~~~~~~~~~~~~~~~~~~~~~ A




% ~~~~~~~~~~~~~~~~~~~~~~~~~~~~~~~~~~~~~~ V frm
\begin{frame}[t,fragile] \frametitle{Side by side: \texttt{\textbackslash lstcoderight}}
  \lstcoderight{codeoutput/primes.py}{Python}{codeoutput/primes.txt}

  \vspace{0.5em}
  \small The code box takes 56\% of the width by default; change it with the first optional
  argument, e.g.\ \verb|\lstcoderight[0.65]{...}{...}{...}|.
\end{frame}
% ~~~~~~~~~~~~~~~~~~~~~~~~~~~~~~~~~~~~~~ A frm



% ~~~~~~~~~~~~~~~~~~~~~~~~~~~~~~~~~~~~~~ V
\subsection[downStep]{Output at below with step}
% ~~~~~~~~~~~~~~~~~~~~~~~~~~~~~~~~~~~~~~ A




% ~~~~~~~~~~~~~~~~~~~~~~~~~~~~~~~~~~~~~~ V frm
\begin{frame}[t,fragile] \frametitle{Output on the next click: \texttt{\textbackslash lstcodedownstep}}
  \lstcodedownstep{codeoutput/grades.py}{Python}{codeoutput/grades.txt}

  \small Slide 1 shows the code; slide 2 adds the output. Turkish letters work in both.
\end{frame}
% ~~~~~~~~~~~~~~~~~~~~~~~~~~~~~~~~~~~~~~ A frm



% ~~~~~~~~~~~~~~~~~~~~~~~~~~~~~~~~~~~~~~ V
\subsection[rightStep]{Output at right with step}
% ~~~~~~~~~~~~~~~~~~~~~~~~~~~~~~~~~~~~~~ A




% ~~~~~~~~~~~~~~~~~~~~~~~~~~~~~~~~~~~~~~ V frm
\begin{frame}[t,fragile] \frametitle{Side by side, output on the next click: \texttt{\textbackslash lstcoderightstep}}
  \lstcoderightstep{codeoutput/fibonacci.py}{Python}{codeoutput/fibonacci.txt}

  \vspace{0.5em}
  \small The output's space is reserved on slide 1, so nothing moves when it appears.
  Text after the macro (like this line) is shown on both slides.
\end{frame}
% ~~~~~~~~~~~~~~~~~~~~~~~~~~~~~~~~~~~~~~ A frm



% ~~~~~~~~~~~~~~~~~~~~~~~~~~~~~~~~~~~~~~ V
\subsection[error]{Errors as output}
% ~~~~~~~~~~~~~~~~~~~~~~~~~~~~~~~~~~~~~~ A




% ~~~~~~~~~~~~~~~~~~~~~~~~~~~~~~~~~~~~~~ V frm
\begin{frame}[t,fragile] \frametitle{Errors are output too}
  \lstcodedown[basicstyle=\ttfamily\scriptsize][basicstyle=\ttfamily\scriptsize]{codeoutput/divide.py}{Python}{codeoutput/divide.txt}

  \small The runner also captures stderr, so tracebacks appear exactly as students see them.
\end{frame}
% ~~~~~~~~~~~~~~~~~~~~~~~~~~~~~~~~~~~~~~ A frm




% ~~~~~~~~~~~~~~~~~~~~~~~~~~~~~~~~~~~~~~ V
\hSection[Overlay]{Output as an Overlay}{}{}
% ~~~~~~~~~~~~~~~~~~~~~~~~~~~~~~~~~~~~~~ A



% ~~~~~~~~~~~~~~~~~~~~~~~~~~~~~~~~~~~~~~ V
\subsection[onTop]{Output on top of the code}
% ~~~~~~~~~~~~~~~~~~~~~~~~~~~~~~~~~~~~~~ A




% ~~~~~~~~~~~~~~~~~~~~~~~~~~~~~~~~~~~~~~ V frm
\begin{frame}[t,fragile] \frametitle{Output on top of the code: \texttt{\textbackslash lstcodeoverlaystep}}
  \lstcodeoverlaystep[][][0.4\textwidth]{8.9cm,2.1cm}{codeoutput/primes.py}{Python}{codeoutput/primes.txt}

  \small \only<1>{Slide 1: only the code. Ask what it prints, then go to the next slide.}%
  \only<2>{Slide 2: the output floats on top, its upper-left corner at \texttt{(8.9cm,2.1cm)}.}
\end{frame}
% ~~~~~~~~~~~~~~~~~~~~~~~~~~~~~~~~~~~~~~ A frm




% ~~~~~~~~~~~~~~~~~~~~~~~~~~~~~~~~~~~~~~ V
\subsection[cover]{Output covers the code}
% ~~~~~~~~~~~~~~~~~~~~~~~~~~~~~~~~~~~~~~ A




% ~~~~~~~~~~~~~~~~~~~~~~~~~~~~~~~~~~~~~~ V frm
\begin{frame}[t,fragile] \frametitle{The overlay can cover part of the code}
  \lstcodeoverlaystep[basicstyle=\ttfamily\small][basicstyle=\ttfamily\scriptsize][0.72\textwidth]%
    {3.5cm,3.6cm}{codeoutput/divide.py}{Python}{codeoutput/divide.txt}

  \vspace{0.3em}
  \small \only<1>{What happens on line 6?}\only<2>{The traceback appears over the code, like a
  terminal window popping up.}
\end{frame}
% ~~~~~~~~~~~~~~~~~~~~~~~~~~~~~~~~~~~~~~ A frm




% ~~~~~~~~~~~~~~~~~~~~~~~~~~~~~~~~~~~~~~ V
\hSection[Part]{Showing a part of a program}{}{}
% ~~~~~~~~~~~~~~~~~~~~~~~~~~~~~~~~~~~~~~ A



% ~~~~~~~~~~~~~~~~~~~~~~~~~~~~~~~~~~~~~~ V
\subsection[allCode]{The code}
% ~~~~~~~~~~~~~~~~~~~~~~~~~~~~~~~~~~~~~~ A




% ~~~~~~~~~~~~~~~~~~~~~~~~~~~~~~~~~~~~~~ V frm
\begin{frame}[t,fragile] \frametitle{A C program}
  \lstcoderight[0.6][basicstyle=\ttfamily\scriptsize][basicstyle=\ttfamily\scriptsize]{codeoutput/factorial.c}{C}{codeoutput/factorial.txt}
\end{frame}
% ~~~~~~~~~~~~~~~~~~~~~~~~~~~~~~~~~~~~~~ A frm



% ~~~~~~~~~~~~~~~~~~~~~~~~~~~~~~~~~~~~~~ V
\subsection[segment]{A segment of the code}
% ~~~~~~~~~~~~~~~~~~~~~~~~~~~~~~~~~~~~~~ A




% ~~~~~~~~~~~~~~~~~~~~~~~~~~~~~~~~~~~~~~ V frm
\begin{frame}[t,fragile] \frametitle{Showing only part of a file}
  \begin{columns}[T]
    \begin{column}{0.56\textwidth}
      \lstcode[firstline=4, lastline=9, firstnumber=4]{codeoutput/factorial.c}{C}
    \end{column}
    \begin{column}{0.42\textwidth}
      \lstoutput[firstline=5, lastline=7]{codeoutput/factorial.txt}
    \end{column}
  \end{columns}

  \vspace{0.5em}


% +++++++++++++++++++++++++++++++++++++++ V lst
%: -Lst.
\begin{lstlisting}[language={[LaTeX]TeX}, basicstyle=\ttfamily\scriptsize, frame=single,
                   backgroundcolor=\color{codebg}, rulecolor=\color{gray!40}]
\lstcode[firstline=4, lastline=9, firstnumber=4]{codeoutput/factorial.c}{C}
\lstoutput[firstline=5, lastline=7]{codeoutput/factorial.txt}
\end{lstlisting}
% +++++++++++++++++++++++++++++++++++++++ A lst


  \small \texttt{firstnumber} keeps the line numbers the same as in the real file.
\end{frame}
% ~~~~~~~~~~~~~~~~~~~~~~~~~~~~~~~~~~~~~~ A frm




% ~~~~~~~~~~~~~~~~~~~~~~~~~~~~~~~~~~~~~~ V
\hSection[Languages]{Some languages}{}{}
% ~~~~~~~~~~~~~~~~~~~~~~~~~~~~~~~~~~~~~~ A






% ~~~~~~~~~~~~~~~~~~~~~~~~~~~~~~~~~~~~~~ V
\subsection[Java]{Java Example}
% ~~~~~~~~~~~~~~~~~~~~~~~~~~~~~~~~~~~~~~ A




% ~~~~~~~~~~~~~~~~~~~~~~~~~~~~~~~~~~~~~~ V frm
\begin{frame}[t,fragile] \frametitle{A Java program}
  \lstcodedown[basicstyle=\ttfamily\scriptsize][basicstyle=\ttfamily\scriptsize]{codeoutput/Hello.java}{Java}{codeoutput/Hello.txt}
\end{frame}
% ~~~~~~~~~~~~~~~~~~~~~~~~~~~~~~~~~~~~~~ A frm




% ~~~~~~~~~~~~~~~~~~~~~~~~~~~~~~~~~~~~~~ V
\subsection[Shell]{Shell Example}
% ~~~~~~~~~~~~~~~~~~~~~~~~~~~~~~~~~~~~~~ A




% ~~~~~~~~~~~~~~~~~~~~~~~~~~~~~~~~~~~~~~ V frm
\begin{frame}[t,fragile] \frametitle{A shell script: Git in action}
  \lstcoderight[0.54][basicstyle=\ttfamily\scriptsize][basicstyle=\ttfamily\scriptsize]{codeoutput/git-demo.sh}{bash}{codeoutput/git-demo.txt}
\end{frame}
% ~~~~~~~~~~~~~~~~~~~~~~~~~~~~~~~~~~~~~~ A frm




% ~~~~~~~~~~~~~~~~~~~~~~~~~~~~~~~~~~~~~~ V
\hSection{Help}{}{}
% ~~~~~~~~~~~~~~~~~~~~~~~~~~~~~~~~~~~~~~ A




% ~~~~~~~~~~~~~~~~~~~~~~~~~~~~~~~~~~~~~~ V
\subsection[Tips]{Tips}
% ~~~~~~~~~~~~~~~~~~~~~~~~~~~~~~~~~~~~~~ A




% ~~~~~~~~~~~~~~~~~~~~~~~~~~~~~~~~~~~~~~ V frm
\begin{frame}[t,fragile] \frametitle{Tips}
  \begin{columns}[T]
    \begin{column}{0.5\textwidth}
      \begin{itemize}\small
        \item \textbf{Too long?} Show an excerpt with \texttt{firstline}/\texttt{lastline},
              or shrink the font with \texttt{basicstyle=\textbackslash ttfamily\textbackslash scriptsize}
        \item \textbf{New language?} Add a line to \texttt{run\_examples.sh}, and use any
              \texttt{listings} language name (\texttt{Python}, \texttt{C}, \texttt{Java},
              \texttt{bash}, \texttt{SQL}, \texttt{Haskell}, \dots)
        \item \textbf{Run one program:} \cmd{sh run\_examples.sh primes.py}
        \item \textbf{Frames} must be \texttt{[fragile]} (they already are in this template)
      \end{itemize}
    \end{column}
    \begin{column}{0.46\textwidth}
      \small \textbf{A missing file does not stop the build.}
      \verb|\lstcode{codeoutput/nothere.py}{Python}| gives:
      \vspace{0.3em}

      \lstcode{codeoutput/nothere.py}{Python}

      \small so you can see at a glance which program still has to be written or run.
    \end{column}
  \end{columns}
\end{frame}
% ~~~~~~~~~~~~~~~~~~~~~~~~~~~~~~~~~~~~~~ A frm




% ~~~~~~~~~~~~~~~~~~~~~~~~~~~~~~~~~~~~~~ V
\subsection[slidegrid]{slidegrid}
% ~~~~~~~~~~~~~~~~~~~~~~~~~~~~~~~~~~~~~~ A




% ~~~~~~~~~~~~~~~~~~~~~~~~~~~~~~~~~~~~~~ V frm
\begin{frame}[t,fragile] \frametitle{Finding $x,y$: \texttt{\textbackslash slidegrid}}
  \slidegrid
  \lstcodeoverlaystep[basicstyle=\ttfamily\scriptsize][][0.4\textwidth]{8.9cm,2.1cm}%
    {codeoutput/primes.py}{Python}{codeoutput/primes.txt}

  \small Put \verb|\slidegrid| in a frame to draw a 1\,cm grid, read off $(x,y)$ for the
  overlay's upper-left corner, then remove it. The slide is 16\,cm $\times$ 9\,cm.
\end{frame}
% ~~~~~~~~~~~~~~~~~~~~~~~~~~~~~~~~~~~~~~ A frm




% ~~~~~~~~~~~~~~~~~~~~~~~~~~~~~~~~~~~~~~ V frm
\begin{frame}[t,fragile]
	\begin{center}
		\vspace{4cm}
		\hRemark{{\Huge {Questions?}}}\\
	\end{center}
\end{frame}
% ~~~~~~~~~~~~~~~~~~~~~~~~~~~~~~~~~~~~~~ A frm



\end{document}
